\section{Evaluation protocol}
\label{sec:setup}
The evaluation links changes in observation to building predictions and task answers. The development camera-ladder diagnostic compares cruise, intermediate, and floor observations for the same 7,238 annotated buildings in 43 ROIs. Its fixed-denominator measure counts unmatched ground-truth buildings as incorrect. Classification accuracy among matched buildings describes a selected subset and is not used as the task endpoint. This diagnostic helps interpret whether the rendered view changes available predicted evidence.

The final online set is separate from the 160-question development set used to select policy rules. Seven frozen policies each answer the same 160 final questions under three generation repeats, yielding $7\times160\times3=3{,}360$ episodes. The repeats vary seeded generation conditions for the same questions; they do not increase the number of independent questions, ROIs, or events. The primary endpoint is terminal question accuracy, with abstentions counted incorrect. Executed dedicated rechecks, within-policy answer transitions, task type, and event are descriptive companion measures.

We audited spatial separation beyond distinct tile IDs. The 43 development ROIs and 44 final ROIs share no tile ID or xBD building UID (0 shared among 7,238 and 4,963 eligible ROI buildings), but 9/1,892 cross-set ROI pairs intersect; the maximum ROI IoU is 0.461 and the maximum fraction of the smaller ROI covered by the other is 0.631. Comparing 129 centered development camera-ladder footprints with all 802 recorded T1 final observation footprints across three repeats, 1,871/103,458 cross-set pairs intersect. The maximum footprint IoU is 0.882, and one smaller footprint lies entirely inside a larger one (overlap fraction 1.000). These measurements identify overlapping geographic windows but do not establish whether both sets contain the same xBD pixels or physical buildings. Different tiles may assign different UIDs to the same physical structure. Thus question and tile disjointness does not establish spatial independence.

We add two controls on the frozen final set. First, the task-majority baseline predicts one answer per task type using a question/ROI-disjoint development set. S1\_LANGUAGE receives question text, choices, and task type; S2\_METADATA additionally receives runtime-visible ROI bounds/tile, initial pose, and, for damage questions, the permitted marked-building reference and coordinate. Both exclude images, detector outputs, annotation labels and polygons, the reference answer, and the spatial true-nearest target. S1 and S2 use the same Qwen checkpoint, answer vocabulary, generation settings, seed schedule, and exact-match scoring as the main answer model, but a modality-appropriate text-only JSON prompt and no hybrid rule path.

Second, a fixed-view matched-call diagnostic is run at every one of the 161 executed T1 dedicated rechecks. Immediately before the live recheck, the fixed branch clones the semantic map and evidence memory, remains at the old pose, reruns perception on the old view, and calls the ordinary hybrid answer path. The changed branch executes the logged action, perceives the new view, and calls that same answer path. The two branches share question, repeat, next-step seed, model, and generation settings; their prompts differ only through the resulting predicted evidence. A strict gate requires a byte-identical fixed-branch marked image and exactly one fixed and changed model call with the shared seed. The diagnostic compares paired candidate answers immediately after the transition, without following both branches to episode termination. The fixed-branch answer is never used by the T1 controller.

Four T1\_TASK contrasts were specified before the final run: against A0\_NO\_RECHECK (internal identifier A0\_HOLD), A1\_RANDOM, A2\_ALWAYS, and A3U\_RAW\_ENTROPY. Comparisons pair the same question and repeat. The 5,000-draw bootstrap resamples ROI clusters, retaining their questions and repeats, and reports Bonferroni-adjusted 98.75\% intervals for the four contrasts. The intervals characterize uncertainty across the 44 observed ROI clusters. The three-event sample does not support inference about transfer to unseen disasters. Task and event breakdowns are descriptive, not additional adjusted tests.

For the matched diagnostic, correctness pairs and answer disagreement are summarized overall and by task, with a 5,000-draw ROI-cluster bootstrap 95\% interval for the changed-minus-fixed difference. Shortcut results report overall and task-wise accuracy; image-aware A0\_NO\_RECHECK and T1 gaps to S1/S2 are paired by question and repeat with the same ROI-cluster bootstrap. These controls assess visual dependence and observation value and are not used for additional policy selection.

The primary cost count records executed dedicated rechecks. Search can also move and render, and policies can make different numbers of later Qwen calls; total movement, rendering, perception, and generation costs are therefore not summarized by the recheck column. A supplementary archive audit reports observation and generation-attempt counts, search actions, state-derived movement on complete trajectories, and episode wall time (Table~\ref{tab:cost-audit}). Episode wall time is measured software runtime, not flight duration. The final protocol identifies the frozen task, model, prompt, and run configurations; the audit scope and planned release are summarized in \ref{app:provenance}. A runnable single-episode scene and scoring replay is described in \ref{app:minimal-reproduction}.
